\documentclass{article}
\usepackage{amsmath}
\title{Motor Solution Bundle Research}
\begin{document}
\maketitle
\section{Scope and State}
A versioned JSON manifest and NPZ array archive bind a scalar-coordinate
MuJoCo model, stateless fixed-gain controller, motor XML and numerical RK4
stage trace. The v2 event is either explicitly absent or a recorded latched force
with contact marked unverified; no contact event is invented. A validated v1 event can be explicitly migrated without adding
unobserved data. This is a restricted research contract, not a complete
Solution Bundle or a clinical or physical acceptance.
\section{Replay and Numerical Evidence}
At each zero-order hold the control law is
\[\tau_{\rm cmd}=\tau_{\rm ff}+\delta\tau_{\rm ff}
+K_p(q_{\rm ref}-q)+K_d(\dot q_{\rm ref}-\dot q).\]
The reconstructed motor plant evaluates post-saturation generalized drive.
The verifier compares all recorded numerical stage arrays exactly and checks
bundle boundary states, sampled references and held forces against the
archive. Final boundary torque repeats the last held value. Stage values are
not continuous-time extrema. All model, XML, controller and archive bytes
are hashed; the manifest, model and runtime digests must be trusted separately.
\section{Limits and Validation}
Only one hold per fixed step, optional existing closure projection, no passive
shaft and no flight are supported. A prescribed force is mapped by the
existing point Jacobian and passed separately from saturated motor input; its
replayed split state is checked, but no on-face contact is established. The directory must remain stable during loading. The
runtime hash covers a declared subset and does not prove independent runtime,
full dependency identity, historical model origin, anatomical limits or
physical actuator provenance. Real matched-swing contact and optimization
remain unqualified. Focused headless tests and reproduction instructions are
listed in the adjacent README; private model inputs and execution receipts stay
outside version control.
\end{document}
